\section{Endpoints conceptuales}

\subsection{Descripcion}

Este documento reservara el catalogo conceptual de endpoints y payloads de alto nivel del backend.

\subsection{Alcance}

La intencion es descriptiva. No reemplaza una especificacion exhaustiva tipo OpenAPI.

\subsection{Definiciones}

\begin{itemize}
  \item \textbf{Endpoint conceptual}: operacion esperada del backend expresada a nivel funcional.
\end{itemize}

\subsection{Reglas}

\begin{itemize}
  \item El documento describira proposito, entrada esperada y salida esperada por operacion.
  \item Los endpoints no podran contradecir reglas del dominio ni limites arquitectonicos.
\end{itemize}

\subsection{Casos borde}

\begin{itemize}
  \item Pendiente de completar despues de cerrar flujos y contratos base.
\end{itemize}

\subsection{Invariantes}

\begin{itemize}
  \item Los contratos conceptuales deben mantener consistencia entre frontend y backend.
\end{itemize}

\subsection{Preguntas abiertas}

\begin{itemize}
  \item Confirmar agrupacion final de endpoints por modulo.
\end{itemize}

